\chapter{Implémentation, limites et exercices}

\section{Organisation du code}
\begin{longtable}{p{5.6cm}p{9.4cm}}
\toprule Fichier & Rôle\\ \midrule
\small\texttt{scripts/\allowbreak build\_vibrance\_model.py} & Lit les paires Lightroom,
échantillonne, construit les tables, impose la monotonie, mesure le facteur
quasi-gris, valide par validation croisée, écrit l'en-tête C++ et la preuve
JSON. Refuse un état absent ou dupliqué ; ne réécrit jamais une preuve.\\
\small\texttt{src/\allowbreak core/\allowbreak VibranceModel.h} & En-tête généré : axes, tables de gain de
chroma, de gain de clarté et de rotation, facteur quasi-gris. Ne se modifie
pas à la main.\\
\small\texttt{src/\allowbreak core/\allowbreak Vibrance.h/\allowbreak .cpp} & Le module : conversions, lecture
interpolée des tables, transformation du pixel, boucle sur l'image.\\
\small\texttt{src/\allowbreak editor/\allowbreak imageeditor.cpp} & Le cadenceur appelle
\texttt{Vibrance::apply} après les réglages tonals, sans traitement propre.\\
\small\texttt{tests/\allowbreak vibrance\_math\_test.cpp} & Tests unitaires du modèle.\\
\small\texttt{tests/\allowbreak analyse\_couleur.py} & Campagne de mesure du rendu réel
d'\IPR{} contre Lightroom, rendu en mémoire.\\
\bottomrule
\end{longtable}

\section{Le traitement d'une image}
La fonction \texttt{Vibrance::apply} reçoit le master Rec.2020 linéaire en
\texttt{float} sur trois canaux. Elle :
\begin{enumerate}
\item ne fait rien si les deux curseurs sont à zéro, si l'image est vide ou
d'un autre type ;
\item parcourt les lignes en parallèle (\texttt{cv::parallel\_for\_}) : chaque
pixel est indépendant, aucune donnée n'est partagée en écriture ;
\item laisse intact tout pixel gris ($R=G=B$), sans aller-retour matriciel qui
introduirait une erreur d'arrondi ;
\item convertit le pixel en sRGB linéaire, applique la transformation, revient
en Rec.2020.
\end{enumerate}

\section{Ce que vérifient les tests}
\begin{itemize}
\item identités exactes : curseurs à zéro, gris, y compris un gris HDR à 4 ;
\item résultats finis pour des valeurs négatives ou supérieures à 1 ;
\item la chroma ne diminue jamais quand la vibrance augmente ;
\item continuité de part et d'autre de zéro ;
\item saturation $-100$ : gris ; saturation seule : clarté et teinte conservées ;
\item comportement Lightroom : gain $+100$ plausible sur une couleur modérée,
plus faible sur une couleur saturée, réponse progressive, quasi-suppression des
couleurs faibles à $-100$ ;
\item rotation : à $+100$ le bleu tourne vers le cyan, à $-50$ dans l'autre sens ;
\item chemin image : un gris Rec.2020 inchangé, une couleur modifiée, un type
non pris en charge ignoré.
\end{itemize}

\section{Limites}
\begin{itemize}
\item Les tables décrivent Lightroom avec le profil Camera Neutral et la
version de processus courante, sur 49 scènes de trois boîtiers Canon. Un autre
profil ou une autre version de Lightroom demande une nouvelle mesure.
\item Les médianes décrivent le comportement typique. Sur certaines images,
comme un grand ciel bleu saturé, Lightroom s'en écarte ; la part de variance
non expliquée, 5\,\% au positif, existe bien.
\item Les couleurs poussées hors du gamut sRGB sont compressées à la sortie,
comme le reste du rendu.
\item Les quasi-gris et les pixels très sombres sont traités par extrapolation
prudente : les mesures y sont rares et bruitées.
\end{itemize}

\section{Exercices}
\paragraph{Exercice 1 -- Logarithme d'un gain.} Une couleur passe d'une chroma
0,05 à 0,08. Calculer $\gamma$. Quel $\gamma$ ramène ensuite la chroma à 0,05 ?

\paragraph{Exercice 2 -- Rotation.} Un pixel a $a=-0{,}02$ et $b=-0{,}09$.
Calculer sa chroma et sa teinte, puis $(a',b')$ après une rotation de
$-5^\circ$. Vérifier que la chroma est conservée.

\paragraph{Exercice 3 -- Interpolation entre valeurs.} Une cellule donne
$\gamma=0{,}35$ à $+50$ et $\gamma=0{,}91$ à $+100$. Quel gain de chroma à
$+75$ ? À $+20$, sachant que $\gamma=0{,}16$ à $+25$ ?

\paragraph{Exercice 4 -- Régression isotone.} Rendre croissante la suite
$(-3{,}69\,;\ -0{,}82\,;\ -0{,}36)$ pour $-100$, $-50$, $-25$, puis la suite
$(0{,}20\,;\ 0{,}15\,;\ 0{,}10)$ pour $+25$, $+50$, $+100$.

\paragraph{Exercice 5 -- Choix d'une variable.} Pourquoi la chroma absolue
explique-t-elle mal le gain positif, alors que la saturation relative
l'explique bien ? Indice : comparer un bleu sombre et un bleu clair d'aspect
également vif.

\section{Corrigés}
\paragraph{Corrigé 1.} $\gamma=\ln(0{,}08/0{,}05)=\ln1{,}6\approx0{,}47$. Le
retour demande $\gamma=\ln(0{,}05/0{,}08)\approx-0{,}47$ : en logarithme, les
gains inverses sont opposés.

\paragraph{Corrigé 2.} $C=\sqrt{0{,}0004+0{,}0081}\approx0{,}0922$ ;
$h=\operatorname{atan2}(-0{,}09,-0{,}02)\approx-102{,}5^\circ$, soit
$257{,}5^\circ$. Avec $\cos(-5^\circ)\approx0{,}9962$ et
$\sin(-5^\circ)\approx-0{,}0872$ :
$a'=-0{,}02\times0{,}9962-(-0{,}09)(-0{,}0872)\approx-0{,}0278$ et
$b'=-0{,}02\times(-0{,}0872)+(-0{,}09)\times0{,}9962\approx-0{,}0879$.
On trouve $C'=\sqrt{0{,}000773+0{,}00773}\approx0{,}0922$ : la chroma est
conservée, et la teinte devient $252{,}5^\circ$, plus proche du cyan.

\paragraph{Corrigé 3.} À $+75$, $t=0{,}5$ : $\gamma=0{,}63$, gain $e^{0{,}63}\approx1{,}88$.
À $+20$, entre 0 et $+25$, $t=0{,}8$ : $\gamma=0{,}128$, gain $\approx1{,}14$.

\paragraph{Corrigé 4.} La première suite est déjà croissante et négative : elle
est inchangée. La seconde viole l'ordre deux fois. On regroupe les deux
premiers termes en leur moyenne $0{,}175$, qui viole encore l'ordre avec
$0{,}10$ ; on regroupe les trois, de moyenne $0{,}15$. Résultat :
$(0{,}15\,;\ 0{,}15\,;\ 0{,}15)$.

\paragraph{Corrigé 5.} À aspect égal, un bleu sombre a une chroma Oklab plus
faible qu'un bleu clair, parce que la chroma croît avec la luminosité. Classer
par chroma mélange donc des couleurs vives sombres et des couleurs ternes
claires, qui ne réagissent pas pareil. La saturation relative, rapport
indépendant de l'échelle, regroupe les couleurs selon leur « vivacité »,
critère que la vibrance utilise.

\section{Glossaire}
\begin{description}
\item[Chroma $C$] distance au gris dans le plan $(a,b)$ d'Oklab, mesure
absolue de l'intensité.
\item[Saturation $s$] $(\max-\min)/\max$ en ProPhoto linéaire, mesure relative.
\item[Teinte $h$] angle de la couleur dans le plan $(a,b)$.
\item[Gain log $\gamma$] logarithme du rapport des chromas après et avant.
\item[Cellule] case de la table, définie par une plage de saturation, de clarté
et de teinte.
\item[Repli] valeur donnée à une cellule trop peu peuplée.
\item[Régression isotone] suite monotone la plus proche des mesures.
\item[Quasi-gris] pixel de chroma inférieure à 0,016, traité par un facteur.
\item[Validation croisée] mesure de l'erreur sur des données exclues de la
construction du modèle.
\end{description}
